% Generated draft: requires mathematical and source review.
\documentclass[11pt]{article}
\usepackage[margin=25mm]{geometry}
\usepackage{fontspec}
\setmainfont{Latin Modern Roman}
\usepackage{amsmath,amssymb,mathtools,unicode-math}
\setmathfont{Latin Modern Math}
\usepackage{enumitem,needspace,xurl,hyperref}
\setlength{\parindent}{0pt}
\setlength{\parskip}{5pt}
\setlength{\emergencystretch}{3em}
\newcommand{\N}{\mathbb{N}}
\newcommand{\Z}{\mathbb{Z}}
\newcommand{\Q}{\mathbb{Q}}
\newcommand{\R}{\mathbb{R}}
\newcommand{\C}{\mathbb{C}}
\newcommand{\F}{\mathbb{F}}
\newcommand{\A}{\mathbb{A}}
\newcommand{\PP}{\mathbb{P}}
\newcommand{\E}{\mathbb{E}}
\newcommand{\eps}{\varepsilon}
\newcommand{\dd}{\,\mathrm{d}}
\begin{document}
\section*{UnsolvedMath ID: 11000046}
\textbf{Source problem code:} catalogue entry
\subsection{Definitions and mathematical statement}
A Kähler group is a group isomorphic to the fundamental group of a compact Kähler manifold. Here a Kähler manifold has a complex structure and a compatible Hermitian metric whose associated real \((1,1)\)-form is closed. Compactness is part of the group-theoretic definition in this source.

Conjecture 4.8 states, unambiguously,
\[
 \forall g\ge3\ \exists X_g\text{ a compact Kähler manifold}:\qquad
 \pi_1(X_g)\cong\operatorname{Mod}_g.
\]
No asphericity of \(X_g\) is required. The source reports that \(\operatorname{Mod}_2\) is not a Kähler group, reference [Ve], via a finite-extension reduction to pure braid groups. Its compressed statement about iterated free-group extensions needs the precise theorem hypotheses; it is not a general obstruction valid for every such extension.

The contextual strategy follows Toledo [To] for nonuniform lattices in \(\mathrm{SU}(n,1)\), \(n\ge3\): begin with a smooth open variety \(V\), seek a projective compactification \(\overline V\) whose boundary has complex codimension at least \(3\), and use a suitable generic surface section avoiding the boundary plus weak Lefschetz to preserve the fundamental group. The extracted text calls the section a ``generic 2-plane'' and loses some bars in \(V\) versus \(\overline V\); its exact section construction and smoothness hypotheses must be recovered. Footnote 6 explicitly calls the presentation a slight cheat and says the actual statement requires orbi-complexes. This caveat is essential when applying the strategy to moduli.

The source says the Deligne--Mumford compactification of \(\mathcal M_g\) has a complex-codimension-one stratum preventing this argument, and that other compactifications are problematic. The exact stratum and compactification notation require checking. The supplied research summary instead interprets the conjecture as merely saying that \(\mathcal M_g\) is Kähler because it is quasiprojective. That is a different statement: a Kähler structure on a noncompact moduli space does not supply the compact manifold required above, and its coarse fundamental group is not automatically its mapping-class-group orbifold fundamental group. The summary's claimed solved interpretation is therefore not substituted for the explicit conjecture.

The source further discusses the Torelli group \(\mathcal I_g\), conventionally the kernel of the homological action \(\operatorname{Mod}_g\to\mathrm{Sp}(2g,\Z)\). It says that for \(g\ge6\) no known Kähler-group constraints rule it out, and cites Hain [Ha3] for a quadratically presented Malcev Lie algebra in that range. These are source-era statements, not a current audit or sufficient conditions for being Kähler. It cites Akita's theorem, called ``Theorem 5.13 above,'' that for \(g\ge7\) its classifying space has no finite-complex homotopy type, ruling out closed aspherical manifold models. It explicitly notes that this does not rule out Kähler groups, mentioning finite Kähler groups as examples. The final sentence refers to a following result [Fa3] not included in the extract; that result is not reconstructed. The summary also reports \(b_1(\operatorname{Mod}_g)=0\) for \(g\ge3\), consistent with the even-first-Betti-number necessary condition for compact Kähler groups; this necessary condition does not establish the conjecture.
\subsection{Short English statement}
The conjecture asks whether every mapping class group of genus at least three is the fundamental group of a compact Kähler manifold. This is distinct from the noncompact moduli space being Kähler. The source's compactification strategy and Torelli discussion are retained with their orbifold, genus and non-asphericity qualifications.
\subsection{Sources}
\begin{itemize}
\item Local UnsolvedMath registry, permanent ID \(11000046\), exact source code \texttt{catalogue entry}, Conjecture 4.8, ``\(\operatorname{Mod}_g\) is Kähler,'' PDF page 40. No individual proposer or proposed year is supplied.
\item Farb, edited volume \emph{Problems on Mapping Class Groups and Related Topics}: \url{https://www.math.uchicago.edu/~farb/papers/mcgbook.pdf}. The volume has 329 explicitly typeset numbered items, repeated local numbers distinguished by PDF page. Extracted context includes footnote 6's orbi-complex warning and a B. Farb footer.
\item Internal references: [Ve] for the genus-two obstruction; Toledo [To] for nonuniform \(\mathrm{SU}(n,1)\) lattices; Hain [Ha3] for the Malcev Lie algebra; Akita, source cross-reference ``Theorem 5.13 above''; and [Fa3] introducing a following unprovided result. Full entries are absent.
\item The research summary names Deligne--Mumford and quasiprojectivity, but explicitly chooses the moduli-space reading despite the source's compact-group definition; its further assertions about mapping-class groups and first Betti number are supplied context, not a verification of this conjecture.
\item Extraction pdf-font-structure, accessed 2026-07-29. Source evidence NEEDS\_REVIEW presents an open problem and calls for a current primary-literature audit. Registry status open and classification OPEN-TRIAGE coexist with a summary beginning PARTIAL and calling its alternative moduli-space interpretation solved; this scope conflict remains explicit.
\item Dataset attribution: Ulam AI / UnsolvedMath Contributors, CC BY 4.0. Source rights NEEDS\_REVIEW: publicly accessible author/publisher copy, redistribution terms requiring review.
\end{itemize}
\end{document}
